\section{Bài toán, mục tiêu và phạm vi}
\subsection{Bài toán nghiên cứu}
Đề tài phân tích chuỗi cảm biến động cơ nhằm ước lượng số cycle còn lại trước khi hỏng. Đầu vào tại cycle $t$ là các quan sát đã có đến $t$ của một động cơ; đầu ra là RUL dự đoán tại cycle đó. Supervised regression (hồi quy có giám sát) là cách mô hình hóa biến mục tiêu dạng số. Mục tiêu báo cáo là dự đoán RUL và giải thích bằng chứng cảnh báo sớm, không chỉ trình bày một thuật toán hồi quy.

Đối với train, đặt $T_i$ là cycle cuối của động cơ $i$ trong lịch sử đầy đủ đến khi hỏng. Nhãn được xác định bởi
\begin{equation}
 y_{i,t}=T_i-t,\qquad 1\leq t\leq T_i.
 \label{eq:rul}
\end{equation}
Ví dụ minh họa: một động cơ kết thúc ở cycle 200 có RUL bằng 30 cycles tại cycle 170 và bằng 0 tại cycle 200. Đây là ví dụ tính nhãn, không phải kết quả dự đoán.

Lịch sử test bị cắt trước khi hỏng nên không được lấy cycle cuối quan sát của test làm $T_i$. Mục tiêu đánh giá cuối là một dự đoán tại cycle cuối quan sát cho mỗi test engine. Bộ nhãn chính thức được tách khỏi quá trình lựa chọn mô hình.

\subsection{Câu hỏi và giới hạn nghiên cứu}
Các câu hỏi chính gồm: dữ liệu có đáp ứng contract không; cảm biến biến đổi thế nào theo vòng đời; đặc trưng từ 30 cycles gần nhất dự đoán RUL được đến mức nào; và kết quả có thể trình bày minh bạch ra sao. Bản hiện tại cung cấp bằng chứng cho câu hỏi chất lượng dữ liệu và một phần EDA, chưa trả lời được câu hỏi độ chính xác mô hình.

FD001 không có chi phí bảo trì, downtime hoặc lịch can thiệp. Do đó, RUL không tự động xác định thời điểm bảo trì tối ưu. Ngưỡng cảnh báo dự kiến 30 cycles chỉ dùng để minh họa cách đọc kết quả. Cycle cũng không được đổi thành giờ hoặc ngày vì chưa có thông tin quy đổi.

\subsection{Liên hệ với yêu cầu môn học}
Theo tài liệu môn học, tiêu chí quy trình xử lý dữ liệu chiếm 40\%, trực quan hóa và phân tích 25\%, chất lượng mã nguồn 20\%, trình bày và phản biện 15\%. Vì vậy, báo cáo dành nội dung riêng cho provenance, contract, ETL, lưu trữ và tái lập, bên cạnh mô hình dự đoán. Các trọng số này là tiêu chí chấm điểm, không phải tỷ lệ hoàn thành hiện tại.

Nguồn yêu cầu là hai tài liệu \textit{Đề cương môn học Phân tích dữ liệu lớn} và \textit{Danh sách đề tài báo cáo cuối kỳ}, được đối chiếu ngày \reportdate. Yêu cầu ghi rõ tiến độ được bổ sung theo trao đổi ngày này. Chưa có mẫu bìa bắt buộc hoặc tên giảng viên được xác nhận trong các nguồn đã đọc; bản hiện tại không tự điền các thông tin đó.

\section{Dữ liệu và kiểm soát chất lượng}
\subsection{Nguồn dữ liệu và lý do lựa chọn}
NASA C-MAPSS mô phỏng các chuỗi vận hành đến khi hỏng của động cơ turbofan. Nguồn dữ liệu được ghi nhận tại NASA Open Data~\cite{nasaCmapssDataset}; tài liệu nền tảng đi kèm bộ dữ liệu là công trình của Saxena và cộng sự~\cite{saxena2008damage}. Mô tả đi kèm file raw xác định FD001 có một điều kiện vận hành và một dạng suy giảm High-Pressure Compressor (HPC, máy nén áp suất cao).

FD001 được chọn vì có lịch sử train đầy đủ đến khi hỏng, test bị cắt và nhãn đánh giá tách riêng, phù hợp với cách xây dựng bài toán RUL. Chỉ sử dụng một tập con giúp kiểm soát phạm vi thực nghiệm và làm rõ từng bước xử lý. Đổi lại, bằng chứng chỉ áp dụng trong điều kiện mô phỏng đang xét; không đủ để khẳng định khả năng khái quát trên động cơ thực tế hoặc điều kiện vận hành khác.

\begin{table}[H]
\centering
\caption{Quy mô FD001 đã ghi nhận trong báo cáo validation.}
\label{tab:data}
\begin{tabularx}{\linewidth}{l r r X}
\toprule Thành phần & Số dòng & Động cơ & Ý nghĩa\\\midrule
Train & 20.631 & 100 & Lịch sử đến khi hỏng\\
Test & 13.096 & 100 & Lịch sử bị cắt trước khi hỏng\\
Test labels & 100 & 100 & RUL tại cycle cuối quan sát\\\bottomrule
\end{tabularx}
\end{table}

Mỗi dòng raw có 26 trường: engine ID, cycle, 3 operational settings và 21 sensors. Reader bổ sung cột \texttt{split}, tạo bảng observations có 27 cột. Khóa là bộ ba \texttt{(split, engine\_id, cycle)}: train engine 1 và test engine 1 không phải cùng một thực thể. Nhãn test chỉ được kiểm tra cấu trúc riêng trong báo cáo integrity; việc này không phải đánh giá mô hình.

\subsection{Các kiểm tra đã triển khai và lý do}
Reader và validator thực hiện kiểm tra số trường, schema, kiểu số hữu hạn, bảng không rỗng, split hợp lệ, khóa duy nhất và cycle liên tục trong từng động cơ. Profile FD001 đối chiếu thêm số dòng, ID động cơ và cycle bắt đầu. Với train targets, validator kiểm tra đủ khóa và đối chiếu Công thức~\ref{eq:rul}.

Các kiểm tra được đặt trước EDA vì dữ liệu sai có thể tạo ra thống kê hợp lệ về mặt tính toán nhưng sai về ý nghĩa. Chẳng hạn, chuỗi cycle $1,2,4$ không được coi là ba bước liên tiếp khi xây window. Hai dòng cùng khóa cũng không được giữ lại như hai phép đo độc lập. Khi phát hiện lỗi, package báo \texttt{ContractError}; không tự điền dữ liệu, xóa dòng hoặc chỉnh raw để làm validation đạt.

\subsection{Tính toàn vẹn và giới hạn của validation}
Báo cáo JSON lưu SHA-256 của input để nhận biết phiên bản dữ liệu đã đọc. Dấu vân tay này chỉ hỗ trợ đối chiếu; tự tính checksum không chứng minh nguồn gốc file là chính xác. Việc kiểm tra row count và cycle continuity cũng không đủ phát hiện mọi thay đổi giá trị sensor.

Đặc biệt, một chuỗi train bị cắt đuôi nhưng vẫn liên tục có thể tạo ra RUL sai nếu bị hiểu nhầm là lịch sử đầy đủ. Vì vậy, giữ raw bất biến và ghi nguồn dữ liệu là điều kiện đi kèm công thức tính nhãn. API test labels được tách riêng nhưng chưa có cơ chế phân quyền filesystem ngăn mọi chương trình đọc file; ranh giới không dùng test để chọn mô hình còn phải được kiểm tra khi xây pipeline mô hình.
